\uuid{aQnm}
\titre{Variables continues}
\auteur{worms}
\organisation{exo7}
\licence{CC BY-NC-SA 4.0}
\datecreate{2017-08-28}
\niveau{L2}
\module{Probabilité et statistique}
\chapitre{Probabilité continue}
\sousChapitre{Densité de probabilité}
\difficulte{}
\theme{Variables continues}
\exo7qcmid{worms.worms.045}

\contenu{
\begin{question}
\qclassification{261.01}
Quelle(s) affirmation(s), parmi les suivantes, sont vraies ?
\begin{answers}
\bad{Si $X$ est de loi continue, alors sa fonction de densité est nécessairement continue sur $\Rr$.}
\good{Si $X$ est de loi continue, alors sa fonction de répartition est nécessairement continue sur $\Rr$.}
\bad{Si $X$ est une variable aléatoire de loi ${\cal E}xp(\lambda)$, alors $2X$ est de loi ${\cal E}xp(2\lambda)$.}
\bad{Si $X$ est une variable aléatoire de loi uniforme sur $[0,1]$, alors $2X$ est de loi de densité $f(x)=\mathbb{I}_{[0,2]}(x)$.}
\end{answers}
\begin{explanations}
C'est la fonction de répartition d'une loi continue qui est continue, mais pas forcément la densité (regardez par exemple les lois uniformes ou exponentielles !). Quand $X\sim {\cal E}xp(\lambda)$, si on avait $2X$ de loi ${\cal E}xp(2\lambda)$, alors $2X$ (qui a comme espérance $2 \mathbb{E}(X)=2/\lambda$) aurait comme espérance $1/(2\lambda)$, ce qui est contradictoire. Quant à la question sur les lois uniformes, il manque un facteur $1/2$ pour que $f$ constitue une fonction de densité.
\end{explanations}
\end{question}
}
